\documentclass[11pt]{article}
\usepackage[T1]{fontenc}
\usepackage[utf8]{inputenc}
\usepackage{amsmath, amssymb}
\usepackage{booktabs}
\usepackage{graphicx}
\usepackage{geometry}
\geometry{margin=2.4cm}
\usepackage{setspace}
\usepackage[bottom]{footmisc}
\usepackage[hidelinks]{hyperref}
\usepackage{threeparttable}
\usepackage{float}
\usepackage[section]{placeins}
\usepackage{caption}
\newcommand{\sym}[1]{$^{#1}$}
\newcommand{\citep}[1]{(\textsc{#1})}
\newcommand{\citet}[1]{\textsc{#1}}
\graphicspath{{/Users/edoardovitella/Documents/work/projects/Paper_PTA/New/Paper/paper_v3/figures/}}

\title{\textbf{Results --- v5 draft}\\[6pt]
\large Environmental Provisions in Chinese PTAs and Export Composition}
\author{Edoardo Vitella}
\date{\today}

\begin{document}
\maketitle

\noindent\rule{\textwidth}{0.4pt}
\footnotesize
Standalone draft of Section~5, following the sequence: full panel $\rightarrow$ collapsed panel
$\rightarrow$ comparison groups $\rightarrow$ sub-indices $\rightarrow$ dynamics, leave-one-out and
the extensive margin. Five subsections, continuous prose, both margins carried together throughout.
Citations render as small-caps placeholders (no bibliography in this file). Nothing in
\texttt{paper\_v3.tex} has been modified.
\normalsize
\noindent\rule{\textwidth}{0.4pt}

\vspace{10pt}
\noindent\emph{Specification recap.} The estimating equation is
\begin{equation}\label{eq:main}
\ln x_{fpdt} = \beta_1\, EP_{dt}\!\times\!g_p + \beta_2\, EP_{dt}\!\times\!b_p
+ \gamma_1\, TD_{dt}\!\times\!g_p + \gamma_2\, TD_{dt}\!\times\!b_p
+ \theta_{fpd} + \theta_{fdt} + \theta_{pt} + \varepsilon_{fpdt}.
\end{equation}

% ============================================================================
\section{Results}\label{sec:results}
% ============================================================================

\subsection{The full panel}\label{sec:fullpanel}

Table~\ref{tab:fullpanel} reports the estimates of equation~\eqref{eq:main} on the full
firm--product--destination--year panel. There are twelve coefficients of interest in it, and they are
easier to read across than down. Columns~(1)--(3) measure environmental content with the World Bank
depth index and columns~(4)--(6) with the TREND provision count; within each block the outcome is log
export value, then log quantity, then log unit value. Panel~A controls for non-environmental World
Bank depth and Panel~B for the independently coded DESTA index \citep{dur2014}. The green and dirty
interactions are estimated jointly, which is why they should be read together: the design asks
whether an environmental chapter tilts a firm's product mix, and a tilt has two ends.

Only one of those ends moves. On the green margin the coefficient is zero in every cell. Export value
gives $-0.0023$ (s.e.\ 0.0039) with the World Bank index and $-0.0001$ (0.0010) with TREND; switching
the depth control to DESTA changes the third decimal and nothing else. The joint $F$-test on all four
composition interactions returns $p = 0.31$ (WB) and $p = 0.71$ (TREND) in Panel~A, and $p = 0.18$ and
$p = 0.26$ in Panel~B. Quantity and unit value add nothing: across the eight green cells in
columns~(2), (3), (5) and~(6) no coefficient comes close to significance, which rules out the one
reading a value regression alone cannot dismiss --- that a volume response and a price response are
offsetting each other inside the value aggregate. Neither channel moves. If Chinese firms shifted
toward environmental goods when an agreement with environmental content entered into force, the shift
is not visible here at any margin.

The dirty coefficient does move, and it is worth being precise about how much. In Panel~A it is
$-0.0044$ on export value with an asymptotic $p$-value of 0.052, and $-0.0009$ with $p = 0.15$ under
TREND. In Panel~B the same two coefficients are $-0.0056$ ($p = 0.039$) and $-0.0014$ ($p = 0.048$).
The sign is stable, the magnitude is stable, and the asymptotic verdict depends on which measure of
non-environmental agreement depth sits in the regression. That dependence has a mechanical
explanation. Net of the fixed effects, environmental depth correlates $0.959$ with World Bank
non-environmental depth and $0.891$ with DESTA depth. The World Bank control absorbs almost the same
variation as the regressor of interest, so little independent variation is left and the standard error
is wide; DESTA overlaps less, leaves more, and the standard error narrows. Nothing about the effect
changes between the panels --- only how much of it the design can see.

This is the point at which asymptotic $p$-values stop being informative. Standard errors clustered at
the destination level are consistent as the number of clusters grows, and this design has 23 treated
destinations sharing roughly 13 distinct agreement profiles, since the ten ASEAN partners enter under a
single agreement. Under those conditions cluster-robust inference over-rejects, sometimes badly. The
row labelled ``Bootstrap $p$'' in every table of this paper reports the wild cluster bootstrap of
\citet{cameron2008}, implemented following \citet{roodman2019} with $B = 9{,}999$ replications: the
model is re-estimated thousands of times with the residuals of each destination flipped in sign at
random, and the observed statistic is placed in the distribution this generates. It is the reference
inference throughout, and the asymptotic figures are kept alongside it only so the gap between the two
remains visible.

The gap is large exactly where it matters. On the green margin the bootstrap changes nothing: $p$
moves from 0.57 to 0.69 (WB) and from 0.92 to 0.93 (TREND), and the null is a null either way. On the
dirty margin under the World Bank control it changes the verdict, taking $p$ from 0.052 to 0.19 and
from 0.15 to 0.18. Under DESTA it does not: 0.039 becomes 0.049 and 0.048 becomes 0.069. So the dirty
coefficient survives robust inference under one depth control and not under the other, and it does so
on both indices at once, which at least means the two codings are telling the same story. The one
place where it is strongest is quantity: under DESTA the World Bank interaction is $-0.0070$ on log
quantity with a bootstrap $p$ of 0.035, larger and better determined than on value, while the unit
value coefficient is $+0.0004$ and nowhere near significance. Whatever this is, it is a shipment
volume, not a price.

The sign, at least, is where the recent literature on China would put it. \citet{yue2024} find that
environmental provisions in Chinese PTAs inhibit exports of energy-consuming products, and
\citet{zhusun2026}, working with customs data that overlap with those used here, report that
environmental provisions reduce dirty exports and raise green ones. On dirty the present estimates
agree in sign and are considerably smaller; on green they do not agree at all. The difference is not
mysterious. Their identification comes from variation across firms and destinations, which is
precisely the variation that $\theta_{fdt}$ removes here. Section~\ref{sec:collapsed} measures how
much of the association that channel accounts for. \citet{zhusun2025} report that environmental
provisions are associated with quality upgrading in Chinese firms' exports; unit value is a coarse
proxy for quality and column~(3) is null, but the outcome and the sample differ enough that this
should not be read as a test of their result.

Before going further it is worth asking how much of Table~\ref{tab:fullpanel} is the specification
rather than the data. Table~\ref{tab:ladder_compact} re-estimates the same interactions with
progressively different sets of absorbed effects. The pattern is instructive. In columns~(1)--(4) the
coefficients are frequently positive and occasionally significant: the World Bank dirty interaction
reaches $+0.0064$ ($p = 0.019$) and the TREND green interaction $+0.0021$ ($p = 0.021$). None of those
four specifications contains a fixed effect with both a destination and a time dimension. Chinese
exports to a partner grow around the entry into force of an agreement for reasons that have nothing to
do with its environmental chapter, and without a destination--time absorber that growth loads onto
whichever interaction is available. Column~(5) adds $\theta_{fdt}$ and the coefficients collapse to the
values of Table~\ref{tab:fullpanel}. The saturated specification is not one choice among five; it is
the only one in the sequence that removes the obvious confounder, and the others are shown to make
clear what it removes.

\input{tabs/tab_fullpanel}

\begin{table}[htbp]
\centering
\footnotesize
\caption{Fixed-effects structure and the composition interactions (full panel, log export value)}
\label{tab:ladder_compact}
\begin{threeparttable}
\begin{tabular}{lccccc}
\toprule
 & (1) & (2) & (3) & (4) & (5) \\
\midrule
$EP^{WB}_{dt} \times g_p$ & 0.0018 & 0.0012 & 0.0060 & 0.0050 & $-$0.0023 \\
 & (0.0050) & (0.0027) & (0.0050) & (0.0037) & (0.0039) \\
\addlinespace[2pt]
$EP^{WB}_{dt} \times b_p$ & $-$0.0037 & $-$0.0015 & 0.0058* & 0.0064** & $-$0.0044* \\
 & (0.0150) & (0.0027) & (0.0031) & (0.0027) & (0.0022) \\
\addlinespace[5pt]
$EP^{TR}_{dt} \times g_p$ & $-$0.0013 & 0.0010 & 0.0021** & 0.0018** & $-$0.0001 \\
 & (0.0010) & (0.0009) & (0.0009) & (0.0008) & (0.0010) \\
\addlinespace[2pt]
$EP^{TR}_{dt} \times b_p$ & $-$0.0049*** & 0.00003 & 0.0002 & 0.0009 & $-$0.0009 \\
 & (0.0015) & (0.0007) & (0.0011) & (0.0008) & (0.0006) \\
\midrule
\quad $\theta_{fpd}$ & \checkmark &  & \checkmark & \checkmark & \checkmark \\
\quad $\theta_{fpt}$ &  & \checkmark & \checkmark &  &  \\
\quad $\theta_{pd}$  &  & \checkmark &  &  &  \\
\quad $\theta_{pt}$  &  &  &  & \checkmark & \checkmark \\
\quad $\theta_{t}$   & \checkmark &  &  &  &  \\
\quad $\theta_{fdt}$ &  &  &  &  & \checkmark \\
\addlinespace[3pt]
Absorbs destination$\times$time? & No & No & No & No & \textbf{Yes} \\
\addlinespace[3pt]
Observations & 29,477,365 & 35,607,036 & 22,927,402 & 29,473,145 & 21,519,511 \\
Clusters ($d$) & 231 & 232 & 231 & 231 & 225 \\
\bottomrule
\end{tabular}
\begin{tablenotes}[flushleft]\footnotesize
\item \textit{Notes:} Each column is a separate estimation of the composition interactions on the full
  firm-level panel with a different set of absorbed fixed effects; the dependent variable is log export
  value throughout. Column~(5) is the reference specification of Table~\ref{tab:fullpanel}, Panel~A,
  column~(1), reproduced here so the sequence terminates at it. Columns~(1)--(4) differ in observation
  count because a less saturated fixed-effects structure removes fewer singleton cells. Hong Kong and
  Macao excluded throughout. All specifications include $TD_{dt} \times g_p$ and $TD_{dt} \times b_p$
  controls (not shown). Standard errors clustered at the destination level in parentheses.
  $^{*}$~$p<0.10$, $^{**}$~$p<0.05$, $^{***}$~$p<0.01$ (asymptotic).
\item Columns~(1)--(4) contain no fixed effect with both a destination and a time dimension, so
  destination-specific shocks that move with the timing of an agreement are not absorbed. Column~(5)
  adds $\theta_{fdt}$. Quantity and unit-value outcomes, and the versions including Hong Kong and Macao,
  are reported in Appendix Tables~\ref{tab:ladder_wb} and~\ref{tab:ladder_trend}.
\item \textit{Source:} Chinese Customs Office; World Bank DTA Database; TREND Database.
\end{tablenotes}
\end{threeparttable}
\end{table}


\FloatBarrier
\subsection{The collapsed panel}\label{sec:collapsed}

Table~\ref{tab:collapsed} repeats the exercise on HS6--destination--year cells weighted by cell size.
Three things change at once and it is worth naming them before reading the numbers. The unit of
observation is a product--destination--year cell rather than a firm--product--destination--year
shipment; the dependent variable is therefore a cell mean of log exports rather than an individual
flow; and the fixed effects become $\theta_{pd} + \theta_{dt} + \theta_{pt}$, since firm identifiers
no longer exist in the collapsed data. The last of these is the substantive change. The firm dimension
of the absorbers is gone, so a shift in \emph{which} firms ship dirty goods to a destination now enters
the coefficient, where in Table~\ref{tab:fullpanel} it was absorbed. The collapsed panel is not a
robustness check on the full panel; it estimates a wider object, and the difference between the two is
informative in its own right.

On green the two panels agree, as they have agreed everywhere so far: $-0.0046$ (bootstrap $p =
0.65$) with the World Bank index, $+0.0018$ ($p = 0.39$) with TREND, and again nothing on quantity or
unit value. On dirty they do not. The World Bank interaction is $-0.0119$ against $-0.0044$ in the
full panel, a factor of 2.7, with an asymptotic $p$-value below 0.001. Since the two specifications
differ by exactly the between-firm channel, the gap measures it: $(0.0119 - 0.0044)/0.0119 = 0.63$, so
roughly three-fifths of the aggregate dirty association is a change in the composition of exporters
rather than reallocation inside a given firm. This is the quantity that reconciles the present null
with \citet{zhusun2026}. Their design keeps that channel; equation~\eqref{eq:main} removes it. Neither
is wrong, and the two results are not in conflict once the estimands are stated.

The collapsed panel is also where the design can be attacked with a third test. The wild cluster
bootstrap already raises the dirty $p$-value from below 0.001 to 0.072, which is a large correction,
but the bootstrap still takes the assignment of treatment as given and only resamples the residuals.
The permutation test asks a blunter question. The environmental profile of each agreement --- its
depth and the years in which it takes effect --- is detached from its destination and reassigned at
random among the 23 treated destinations; the model is re-estimated; and the exercise is repeated a
thousand times. The reported $p$-value is the share of these random reassignments that produce a
coefficient at least as large in absolute value as the observed one. It is expensive, which is why it
is run on export value only, and it is discrete, because thirteen distinct profiles do not generate a
smooth distribution. But it answers the question a reader of this table actually has: would this
number have appeared if we had simply relabelled which of these destinations signed environmental
chapters?

For the dirty coefficient the answer is that it would, 27.8\% of the time. For green it is 59.7\%.
Under the DESTA control the two figures are 14.6\% and 48.0\%. This is the strongest evidence in the
paper against a causal reading of the dirty coefficient, and it is strong precisely because it is
indifferent to the choice that Section~\ref{sec:fullpanel} could not resolve: the permutation
distribution is built by shuffling treatment, not by changing the control, so it gives the same
verdict under both depth measures. The dirty association is real in the sense that it is in the data;
it is not distinguishable from what random relabelling of these same destinations produces.

\input{tabs/tab_collapsed}

\FloatBarrier
\subsection{Changing the comparison group}\label{sec:comparison}

Everything so far compares green and dirty products against neutral products, and agreement
destinations against non-agreement destinations. Both comparisons can be questioned, and the next
three tables replace them one at a time.

Neutral products may be poor counterfactuals for green ones if the two sit in different regions of the
product space and are exposed to different shocks. Table~\ref{tab:cprodhs4} therefore restricts the
sample to HS6 lines with a comparable match at the four-digit level, so that a solar panel is compared
with its neighbours inside the same heading rather than with the entire neutral universe. The green
coefficient becomes $-0.0009$ with the World Bank index and $+0.0005$ with TREND, both with bootstrap
$p$-values above 0.72. The dirty coefficient does something more interesting: it changes sign, to
$+0.0066$ and $+0.0033$, and remains insignificant. Narrowing the product comparison does not shrink
the dirty estimate towards zero so much as scramble it, which is what one expects of a coefficient
that is not identifying a stable object.

Destinations raise the symmetric concern. China does not sign agreements at random, and if partners
are on different growth paths than non-partners the comparison is contaminated. Table~\ref{tab:cem}
matches destinations by coarsened exact matching on log GDP per capita, growth and the MFN tariff in
2000, which leaves 57 destination clusters that are comparable on pre-treatment observables. Both
margins survive intact: green is $-0.0023$ (WB) and $-0.0006$ (TREND) with bootstrap $p$-values of
0.61, and dirty is $-0.0041$ and $-0.0011$ with bootstrap $p$-values of 0.25 and 0.20. Under the DESTA
control the dirty coefficients are $-0.0043$ and $-0.0013$, with bootstrap $p$-values of 0.12 and 0.13
--- weaker than in the full sample, which is the opposite of what a genuine effect concentrated among
comparable partners would produce.

The third table removes the untreated comparison group altogether. In Table~\ref{tab:deepshallow} the
sample is restricted to the 23 PTA partners, so that identification comes from variation in
environmental depth among destinations that all have an agreement in force: a shallow environmental
chapter is the counterfactual for a deep one. This is the most demanding of the three designs and the
one with the fewest clusters. Every coefficient is small and none is significant --- green $-0.0022$
and $-0.0004$, dirty $-0.0034$ and $-0.0006$ under the World Bank control, with bootstrap $p$-values
between 0.50 and 0.75.

Read together, the three tables say two different things about the two margins. The green null does
not depend on the comparison group: across the six estimates in this section it stays between $-0.0023$
and $+0.0005$ and its bootstrap $p$-value never falls below 0.58. An artifact of the control group
would move when the control group is replaced three times over; this does not move. The dirty
coefficient, by contrast, survives none of the three replacements. It keeps its sign under matching
and among partners, loses it under the within-heading product comparison, and is insignificant in all
of them. Combined with the permutation result of Section~\ref{sec:collapsed}, the reading that fits
the evidence is that the dirty coefficient is a feature of the specification in which it appears
rather than a composition effect that would show up wherever one looked.

\input{tabs/tab_cprodhs4}

\input{tabs/tab_cem}

\input{tabs/tab_deepshallow}

\FloatBarrier
\subsection{Does the content of the agreement matter?}\label{sec:subindices}

One reading of the null so far is that it is an artifact of aggregation. The indices used to this point
add up every environmental provision in an agreement regardless of what it does, and most of what they
add up is cooperation language, institutional machinery and general environmental statements --- text
with no channel through which it could move a product mix. Only a small minority of provisions carry a
trade mechanism: tariff reductions on environmental goods, non-regression clauses that prevent a
partner from relaxing standards to attract trade, binding market-access commitments. If the few
clauses that could bite are diluted by the many that cannot, a null on the aggregate index would say
nothing about whether environmental content matters.

The dilution is severe. The World Bank index checks 150 provision slots across the treated
destinations, and 8 of them --- 5.3\% --- fall in sub-indices with a trade mechanism. TREND, which
codes at finer granularity, gives 4 of 437, or 0.9\%. At the level of the agreement, a binding
environmental obligation appears in exactly one of the fourteen Chinese PTAs in the database, the 2015
agreement with Korea. This is worth stating plainly because it bounds what any heterogeneity analysis
on this sample can deliver: the two World Bank sub-indices with a trade mechanism are non-zero in only
three destination--years and always in the same proportion, which makes them perfectly collinear and
impossible to separate. \citet{brandi2020} report a maximum variance inflation factor of 4.6 across
provision types on a panel of 680 agreements; with thirteen distinct profiles there is no equivalent
here.

What can still be done is to replace the aggregate index with one sub-index at a time and re-estimate.
Table~\ref{tab:subindex_fullpanel} does this on the full panel for seven sub-indices, two from the
World Bank coding and five from TREND. The coefficients are not partial effects holding other clause
types fixed; each absorbs whatever correlated variation moves with it, so the table is a set of
separate comparisons rather than a decomposition. It should be read for pattern, not for magnitude.

The pattern is the same one as before. On the green margin nothing registers: the two trade-mechanism
sub-indices give $+0.0085$ (World Bank Green Liberalization, $p = 0.86$) and $+0.0158$ (TREND Green
Market Access, $p = 0.38$), and the remaining five are equally silent. On the dirty margin two
sub-indices are marginally negative under asymptotic inference --- enforcement and dispute settlement
attached to the environmental chapter, at $-0.0243$ ($p = 0.040$) in the World Bank coding and
$-0.0107$ ($p = 0.081$) in TREND --- and this is the same asymptotics-only signal documented in
Section~\ref{sec:fullpanel}, on a smaller and more collinear regressor. The wild cluster bootstrap was
not run on the full-panel sub-indices, so these two figures are asymptotic and should be discounted
accordingly.\footnote{On the collapsed panel, where the bootstrap is available for the sub-indices,
the only coefficient that survives it is the TREND regulatory-space index, at $+0.024$ on green
($p = 0.045$) and $+0.023$ on dirty ($p = 0.021$). Both margins move together and by the same amount,
so the green--dirty differential is $+0.0017$ with $p = 0.80$: this is not a shift in composition
toward cleaner goods but a proportional rise in both categories against the neutral baseline. The
regulatory-space index also accounts for 71.5\% of the TREND count and correlates 0.90 with total
depth, which makes it a poor placebo rather than a clean channel. Appendix Table~\ref{tab:subindex}
reports the full set.}

The conclusion is a negative one and it is the honest one: this sample cannot tell whether the null
reflects the absence of an effect or the absence of provisions capable of producing one. What it can
establish is that the two are observationally equivalent here, and that a null on the aggregate index
is exactly what \citet{brandi2020}'s own heterogeneity would predict for a set of agreements composed
almost entirely of cooperation clauses.

\input{tabs/tab_subindex_fullpanel}

\FloatBarrier
\subsection{Dynamics, influence, and the extensive margin}\label{sec:closing}

Three final checks, each aimed at a different way the preceding results could be wrong.

The first is timing. Figure~\ref{fig:es} plots the green and dirty gap relative to neutral products
around entry into force, with never-treated destinations as controls. The pre-treatment coefficients
are flat on both margins, which is the standard falsification of the identifying assumption. It should
be read as exactly that and no more: parallel trends is a statement about the counterfactual path
after treatment, and a flat pre-trend rules out the most detectable violation without establishing
that the assumption holds thereafter. A mild negative drift appears for green products in the late
post-period, but it is identified only by the four earliest cohorts and it disappears under the
\citet{sun2021} interaction-weighted estimator, whose aggregated treatment effect is $-0.042$
($p = 0.27$) on the green gap and $+0.073$ ($p = 0.28$) on the dirty gap. The full event-study
coefficients, the Sun--Abraham estimates, and the destination-specific trend specifications that
address selection on trends are reported in the appendix.

The second is influence. A coefficient estimated from 23 treated destinations can be the product of
one of them, and Table~\ref{tab:loo_new} re-estimates the dirty interaction dropping each in turn. The
familiar concern does not materialise: the point estimate stays between $-0.0097$ and $-0.0133$ in all
23 samples, and dropping large partners such as India or Pakistan moves it by one to four per cent.
What is fragile is not the estimate but its precision. Dropping Australia moves the coefficient by
13\%, to $-0.0103$, while nearly tripling its standard error from 0.0030 to 0.0087; dropping Korea
doubles the standard error and gives $-0.0097$. Australia and Korea are not outliers pulling the
estimate --- they are supplying the variation that identifies it, and without them the design cannot
determine much of anything. Which of the two is pivotal, moreover, depends on the depth control:
under DESTA it is Korea rather than Australia. A coefficient whose decisive observation is selected by
a modelling choice unrelated to environmental content is behaving the way the permutation test
already suggested it was.

The third is the margin itself. Everything above conditions on positive flows, so an agreement that
created green exports to a market previously served by none would leave no trace in it.
Table~\ref{tab:ppml} estimates the same interactions by Poisson pseudo-maximum likelihood on an
HS6--destination--year grid completed with zeros \citep{santossilva2006}, of which 7.9 million cells
of 8.2 survive the estimator's removal of separated and singleton observations. The green interaction
is $+0.0015$ ($p = 0.74$) with the World Bank index and $+0.0001$ ($p = 0.95$) with TREND: no trade
creation. The dirty interaction is $-0.030$ ($p = 0.057$) and $+0.003$ ($p = 0.55$), marginal at best
and, once again, only under one of the two codings. One limitation deserves stating: this grid detects
a product newly exported to a destination, not a new firm beginning to export a product that other
Chinese firms already ship there. The extensive margin is probed, not exhausted.

\begin{figure}[htbp]
\centering
\caption{Event study: green and dirty vs.\ neutral products around PTA entry into force}
\label{fig:es}
\vspace{6pt}
\includegraphics[width=\textwidth]{eventstudy_collapsed_v3.png}
\begin{minipage}{\textwidth}\footnotesize
\emph{Notes:} Collapsed panel; FE $pd + dt + pt$; destination-clustered 95\% confidence bands;
$t = -1$ is the reference period. Coefficients are tabulated in Appendix
Table~\ref{tab:eventstudy_twfe}. Source: Chinese Customs Office; World Bank DTA Database.
\end{minipage}
\end{figure}

\input{tabs/tab_loo_new}

\input{tabs/tab_ppml}

\FloatBarrier
\vspace{14pt}
\noindent\rule{\textwidth}{0.4pt}
\subsection*{Note for the author}
\footnotesize
\textbf{In appendix, cited in the text above:} full saturation ladder for all three outcomes
(\ref{tab:ladder_wb}, \ref{tab:ladder_trend}); Sun--Abraham event study (\ref{tab:sunab}); tabulated
TWFE event study (\ref{tab:eventstudy_twfe}); destination-specific trends and pre-trend tests
(\ref{tab:trends}); Hong Kong and Macao variants of both main tables
(\ref{tab:fullpanel_inclHKMO}, \ref{tab:collapsed_inclHKMO}); complete bootstrap and permutation
tables (\ref{tab:wcb}, \ref{tab:perm}); the dirty-coefficient matrix (\ref{tab:matrice});
full-panel robustness battery (\ref{tab:robust-full}); APEC versus OECD product definition
(\ref{tab:apec}); CO$_2$ intensity (\ref{tab:co2}); collapsed-panel sub-indices
(\ref{tab:subindex}); minimum detectable effects (\ref{tab:mde}) and the Brandi benchmark
(\ref{tab:brandi}). C-overlap is dropped, as agreed.

\medskip
\textbf{Three things to decide.} (i)~The bounding argument --- minimum detectable effects and the
comparison with \citet{brandi2020} --- has no home in this version. It is what turns ``we find no
effect'' into ``we can rule out effects above 3\% per provision''. It would fit as two paragraphs at
the end of Section~\ref{sec:fullpanel} or as the opening of the conclusion. (ii)~The wild cluster
bootstrap was never run on the full-panel sub-indices; Section~\ref{sec:subindices} currently reports
asymptotic $p$-values and says so. (iii)~The compact ladder in Table~\ref{tab:ladder_compact} is new;
if you would rather keep the section leaner it moves to the appendix and the argument survives as one
sentence.
\normalsize

\newpage
\appendix
\section*{Appendix placeholders (labels only)}
\footnotesize
\begin{itemize}\itemsep1pt
\item Saturation ladder \phantomsection\label{tab:ladder_wb}\label{tab:ladder_trend}
\item Sun--Abraham \phantomsection\label{tab:sunab}
\item Event study TWFE \phantomsection\label{tab:eventstudy_twfe}
\item Destination trends \phantomsection\label{tab:trends}
\item Full panel incl.\ HK/Macao \phantomsection\label{tab:fullpanel_inclHKMO}
\item Collapsed incl.\ HK/Macao \phantomsection\label{tab:collapsed_inclHKMO}
\item Bootstrap \phantomsection\label{tab:wcb}
\item Permutation \phantomsection\label{tab:perm}
\item Dirty matrix \phantomsection\label{tab:matrice}
\item Robustness battery \phantomsection\label{tab:robust-full}
\item APEC vs OECD \phantomsection\label{tab:apec}
\item CO$_2$ intensity \phantomsection\label{tab:co2}
\item Collapsed sub-indices \phantomsection\label{tab:subindex}
\item Minimum detectable effects \phantomsection\label{tab:mde}
\item Brandi benchmark \phantomsection\label{tab:brandi}
\end{itemize}
\normalsize

\end{document}
